% Generated draft: requires mathematical and source review.
\documentclass[11pt]{article}
\usepackage[margin=25mm]{geometry}
\usepackage{fontspec}
\setmainfont{Latin Modern Roman}
\usepackage{amsmath,amssymb,mathtools,unicode-math}
\setmathfont{Latin Modern Math}
\usepackage{enumitem,needspace,xurl,hyperref}
\setlength{\parindent}{0pt}
\setlength{\parskip}{5pt}
\setlength{\emergencystretch}{3em}
\newcommand{\N}{\mathbb{N}}
\newcommand{\Z}{\mathbb{Z}}
\newcommand{\Q}{\mathbb{Q}}
\newcommand{\R}{\mathbb{R}}
\newcommand{\C}{\mathbb{C}}
\newcommand{\F}{\mathbb{F}}
\newcommand{\A}{\mathbb{A}}
\newcommand{\PP}{\mathbb{P}}
\newcommand{\E}{\mathbb{E}}
\newcommand{\eps}{\varepsilon}
\newcommand{\dd}{\,\mathrm{d}}
\begin{document}
\section*{UnsolvedMath ID: 800086}
\textbf{Source problem code:} catalogue entry
\subsection{Definitions and mathematical statement}
Let \(M\) be a closed hyperbolic \(3\)-manifold with first Betti number \(b_1(M)=\dim_{\R}H^1(M;\R)\geq 1\). An integral real cohomology class means an element of
\[
\operatorname{im}\bigl(H^2(M;\Z)\longrightarrow H^2(M;\R)\bigr).
\]
Write \(x_M\) for the Thurston seminorm on \(H_2(M;\R)\). In the conventional orientable setting, for an integral homology class \(z\), it is given by
\[
x_M(z)=\min_{[S]=z}\sum_i\max\{0,-\chi(S_i)\},
\]
where \(S=\bigsqcup_i S_i\subset M\) ranges over embedded closed oriented surfaces representing \(z\); it extends homogeneously and continuously to real homology. Define its dual, allowing the value \(+\infty\), by
\[
x_M^*(a)=\inf\left\{C\geq 0:\ |\langle a,z\rangle|\leq Cx_M(z)\text{ for every }z\in H_2(M;\R)\right\},
\]
with \(\inf\varnothing=+\infty\).

Let \(\mathrm{Homeo}_+(S^1)\) be the group of orientation-preserving homeomorphisms of the circle. For a representation \(\rho:\pi_1(M)\to\mathrm{Homeo}_+(S^1)\), let \(e_{\Z}(\rho)\in H^2(M;\Z)\) denote the Euler class of its associated oriented flat circle bundle \(\widetilde M\times_{\rho}S^1\to M\), and let \(e_{\R}(\rho)\) be its image in real cohomology.

The supplied question asks whether, for every such \(M\),
\[
\forall a\in\operatorname{im}\bigl(H^2(M;\Z)\to H^2(M;\R)\bigr),\qquad
x_M^*(a)=1\ \Longrightarrow\ \exists\rho:\pi_1(M)\to\mathrm{Homeo}_+(S^1)
\quad e_{\R}(\rho)=a.
\]
The source does not explicitly assume that \(M\) is orientable. The surface description above specifies the usual orientable interpretation; whether nonorientable manifolds are intended, and which Thurston-norm convention applies to them, requires clarification. Since the source places \(a\) in real cohomology, the displayed conclusion concerns the real image of the Euler class. It does not prescribe an integral lift of \(a\), including its torsion component.
\subsection{Short English statement}
For a closed hyperbolic \(3\)-manifold with positive first Betti number, must every integral real second cohomology class of dual Thurston norm exactly one be the real Euler class of an orientation-preserving circle action of its fundamental group? The intended orientability assumption and any requirement concerning a particular integral lift remain unspecified.
\subsection{Sources}
\begin{itemize}
\item Local UnsolvedMath registry: supplied source record, permanent ID \(800086\), source problem code \texttt{catalogue entry}, titled \emph{Foliations and Diffeomorphism Groups --- Question 27}. The registry attributes the problem to Nariman and Yazdi (2025) and labels it \emph{open}; this label is supplied curation, not an independently verified mathematical or literature assessment.
\item Nariman and Yazdi, \emph{Problem List on Foliations and Diffeomorphism Groups} (2025), Question 27. Supplied source URL: \url{https://arxiv.org/abs/2503.18274}. The supplied provenance records access on 2026-10-16 and extraction by \texttt{arxiv\_latex}; these are source-record provenance claims, not a claim of access or verification for this draft.
\end{itemize}
\end{document}
